% GENERATED FILE. Edit canonical lessons, then run npm run generate:books.
% canonical-source-hash: fnv1a64:9f6d76246fabe5ba
% canonical-lessons: IT-C409-forare, IT-C409-foratura, IT-C409-pneumatico, IT-C409-cric, IT-C409-gonfiare

\chapter{Una gomma a terra}
\label{ch:it-una-gomma-a-terra}
\hlchaptermodality{409}

\begin{chapteropening}
\textbf{By the end of this chapter:} \emph{I can deal with a flat tyre: getting a puncture, the tyre, the jack and pumping it up.}

Complete the last lesson of chapter 409: I can deal with a flat tyre: getting a puncture, the tyre, the jack and pumping it up.
\end{chapteropening}

\section[forare]{forare — to get a puncture; to pierce}

\label{lesson:IT-C409-forare}



\begin{quote}
Before the new one: say the Italian for the tank, then the Italian for to refuel.
\end{quote}

\subsection*{You'll want to know: forare}
\textbf{forare} — \textquotedblleft{}to get a puncture; to pierce\textquotedblright{}.

\begin{quote}
Abbiamo forato in autostrada.
\end{quote}

\begin{quote}
\emph{We got a puncture on the motorway.}
\end{quote}

\emph{Ho forato} — I have a flat tyre; \emph{forare un foglio} — to punch a hole in a sheet.

\begin{cousinweb}
\emph{forare} is inherited from Latin \emph{forāre} 'to bore, to pierce' (EtymDB).
\end{cousinweb}

\begin{grammarlens}[title={What you've built}]
One more word for telling, explaining and describing at B1.
\end{grammarlens}

\subsection*{Your turn}
\begin{itemize}
\raggedright
  \item \emph{Say it:} \emph{forare}
  \item \emph{Say it:} \emph{forare}, once more
  \item \emph{Say it:} a sentence of your own with \emph{forare}
  \item \emph{Recall:} say the Italian for the tank, then the Italian for to refuel, then say \emph{forare} again
\end{itemize}

\begin{tcolorbox}[breakable,colback=teal!4,colframe=teal!35!black,title={Before you move on}]
What does forare mean? (\textquotedblleft{}To get a puncture; to pierce\textquotedblright{}.) Say it once more.
\end{tcolorbox}
\section[la foratura]{la foratura — the puncture}

\label{lesson:IT-C409-foratura}



\begin{quote}
Before the new one: say the Italian for to refuel, then the Italian for to get a puncture.
\end{quote}

\subsection*{You'll want to know: la foratura}
\textbf{la foratura} — \textquotedblleft{}the puncture\textquotedblright{}.

\begin{quote}
Per una foratura siamo arrivati tardi.
\end{quote}

\begin{quote}
\emph{Because of a puncture we arrived late.}
\end{quote}

\emph{Riparare una foratura} — to fix a puncture.

\begin{cousinweb}
\emph{foratura} is built on \emph{forare} 'to pierce', which is inherited from Latin \emph{forāre} 'to bore' (EtymDB).
\end{cousinweb}

\begin{grammarlens}[title={What you've built}]
One more word for telling, explaining and describing at B1.
\end{grammarlens}

\subsection*{Your turn}
\begin{itemize}
\raggedright
  \item \emph{Say it:} \emph{la foratura}
  \item \emph{Say it:} \emph{la foratura}, once more
  \item \emph{Say it:} a sentence of your own with \emph{la foratura}
  \item \emph{Recall:} say the Italian for to refuel, then the Italian for to get a puncture, then say \emph{la foratura} again
\end{itemize}

\begin{tcolorbox}[breakable,colback=teal!4,colframe=teal!35!black,title={Before you move on}]
What does la foratura mean? (\textquotedblleft{}The puncture\textquotedblright{}.) Say it once more.
\end{tcolorbox}
\section[lo pneumatico]{lo pneumatico — the tyre}

\label{lesson:IT-C409-pneumatico}



\begin{quote}
Before the new one: say the Italian for to get a puncture, then the Italian for the puncture.
\end{quote}

\subsection*{You'll want to know: lo pneumatico}
\textbf{lo pneumatico} — \textquotedblleft{}the tyre\textquotedblright{}.

\begin{quote}
Gli pneumatici da neve sono obbligatori qui.
\end{quote}

\begin{quote}
\emph{Snow tyres are required here.}
\end{quote}

In speech, \emph{la gomma}. The article is \emph{lo}, plural \emph{gli pneumatici}, though \emph{il pneumatico} is heard too.

\begin{cousinweb}
\emph{pneumatico} is a learned borrowing from Latin \emph{pneumaticus}, from Greek \emph{pneumatikós} 'of wind or air', built on \emph{pneûma} 'air, breath' (EtymDB): a tyre full of air.
\end{cousinweb}

\begin{grammarlens}[title={What you've built}]
One more word for telling, explaining and describing at B1.
\end{grammarlens}

\subsection*{Your turn}
\begin{itemize}
\raggedright
  \item \emph{Say it:} \emph{lo pneumatico}
  \item \emph{Say it:} \emph{lo pneumatico}, once more
  \item \emph{Say it:} a sentence of your own with \emph{lo pneumatico}
  \item \emph{Recall:} say the Italian for to get a puncture, then the Italian for the puncture, then say \emph{lo pneumatico} again
\end{itemize}

\begin{tcolorbox}[breakable,colback=teal!4,colframe=teal!35!black,title={Before you move on}]
What does lo pneumatico mean? (\textquotedblleft{}The tyre\textquotedblright{}.) Say it once more.
\end{tcolorbox}
\section[il cric]{il cric — the jack (for a car)}

\label{lesson:IT-C409-cric}



\begin{quote}
Before the new one: say the Italian for the puncture, then the Italian for the tyre.
\end{quote}

\subsection*{You'll want to know: il cric}
\textbf{il cric} — \textquotedblleft{}the jack (for a car)\textquotedblright{}.

\begin{quote}
Senza il cric non possiamo cambiare la ruota.
\end{quote}

\begin{quote}
\emph{Without the jack we cannot change the wheel.}
\end{quote}

It sounds like the clicking it makes. Plural unchanged: \emph{i cric}.

\begin{cousinweb}
\emph{cric} came from French \emph{cric}, a word that seems to imitate the click of the mechanism; the deeper origin is not settled.
\end{cousinweb}

\begin{grammarlens}[title={What you've built}]
One more word for telling, explaining and describing at B1.
\end{grammarlens}

\subsection*{Your turn}
\begin{itemize}
\raggedright
  \item \emph{Say it:} \emph{il cric}
  \item \emph{Say it:} \emph{il cric}, once more
  \item \emph{Say it:} a sentence of your own with \emph{il cric}
  \item \emph{Recall:} say the Italian for the puncture, then the Italian for the tyre, then say \emph{il cric} again
\end{itemize}

\begin{tcolorbox}[breakable,colback=teal!4,colframe=teal!35!black,title={Before you move on}]
What does il cric mean? (\textquotedblleft{}The jack (for a car)\textquotedblright{}.) Say it once more.
\end{tcolorbox}
\section[gonfiare]{gonfiare — to pump up, to inflate}

\label{lesson:IT-C409-gonfiare}



\begin{quote}
Before the new one: say the Italian for the tyre, then the Italian for the jack.
\end{quote}

\subsection*{You'll want to know: gonfiare}
\textbf{gonfiare} — \textquotedblleft{}to pump up, to inflate\textquotedblright{}.

\begin{quote}
Dal benzinaio puoi gonfiare le gomme.
\end{quote}

\begin{quote}
\emph{At the petrol station you can pump up the tyres.}
\end{quote}

\emph{Gonfiare un palloncino} — to blow up a balloon.

\begin{cousinweb}
\emph{gonfiare} is inherited from Latin \emph{conflāre} 'to blow together, to kindle' (EtymDB), built on \emph{flāre} 'to blow'.
\end{cousinweb}

\begin{grammarlens}[title={What you've built}]
One more word for telling, explaining and describing at B1.
\end{grammarlens}

\subsection*{Your turn}
\begin{itemize}
\raggedright
  \item \emph{Say it:} \emph{gonfiare}
  \item \emph{Say it:} \emph{gonfiare}, once more
  \item \emph{Say it:} a sentence of your own with \emph{gonfiare}
  \item \emph{Recall:} say the Italian for the tyre, then the Italian for the jack, then say \emph{gonfiare} again
\end{itemize}

\begin{tcolorbox}[breakable,colback=teal!4,colframe=teal!35!black,title={Before you move on}]
What does gonfiare mean? (\textquotedblleft{}To pump up, to inflate\textquotedblright{}.) Say it once more.
\end{tcolorbox}
